% !TEX root = ../main.tex
\chapter{Linear Transformations and Change of Basis}
\label{ch:10}

For nine chapters a matrix has been an array of numbers and $Ax$ has been a combination of its columns. That is the right working picture, but it hides where matrices come from. A matrix is the \emph{shadow} of something more basic and coordinate-free: a \emph{linear transformation}, a rule $T$ that takes vectors to vectors and respects addition and scaling. The transformation knows nothing about rows, columns, or numbers. It just says ``project onto this line'' or ``rotate by $45^\circ$'' or ``take the derivative.'' The matrix appears only after we commit to a \emph{basis}, which installs coordinates and turns the abstract rule into multiplication by an array.

Two ideas follow from that one shift in viewpoint, and they organize the chapter. First, because the matrix depends on the basis we chose, the \emph{same} transformation has many matrices --- one per basis --- and they are all linked by the relation $B = M^{-1}AM$. That is the precise meaning of \emph{similar matrices}, and it is why a clever basis (eigenvectors, Fourier vectors, wavelets) can replace a complicated matrix by a simple one. We will see this pay off in image compression. Second, we return to the matrices that have no honest inverse --- rectangular ones, and square ones that are singular --- and ask for the best substitute. Full column rank buys a \emph{left} inverse; full row rank buys a \emph{right} inverse; and in the hardest case, where the rank is short on both sides, the singular value decomposition hands us the \emph{pseudoinverse} $A^{+}$, which inverts $A$ on the part of space where inversion is possible and gives the shortest least-squares solution everywhere else. The four fundamental subspaces, met in Chapter~\ref{ch:4}, run through all of it.

\section{Linear Transformations Without Coordinates}

In older courses linear algebra began here, with transformations, and only later brought in matrices. The geometric approach has a real advantage: it avoids coordinates until computation forces them on us, and it keeps the eye on what the map \emph{does} rather than on a grid of numbers.

A \emph{transformation} $T$ is just a rule that assigns to each input vector $v$ an output vector $T(v)$. We write $T:\RR^n\to\RR^m$ when inputs live in $\RR^n$ and outputs in $\RR^m$. The transformations worth studying are the ones that interact properly with the two operations that define a vector space --- adding and scaling.

\defn{Linear Transformation}{A transformation $T:\RR^n\to\RR^m$ is \textbf{linear} if it respects vector addition and scalar multiplication:
\[
   T(v+w) = T(v) + T(w)
   \qquad\text{and}\qquad
   T(cv) = c\,T(v)
\]
for all vectors $v,w$ and all scalars $c$. The two conditions combine into the single requirement
\[
   T(cv + dw) = c\,T(v) + d\,T(w),
\]
which says that $T$ carries every linear combination of inputs to the same linear combination of outputs.}

Notice what the definition does \emph{not} mention: no rows, no columns, no numbers. It is a statement about combinations, and it is exactly the property that lets a transformation be pinned down by its action on a few vectors. One small consequence is worth recording: a linear transformation must fix the origin,
\[
   T(0) = T(0\cdot v) = 0\cdot T(v) = 0 .
\]
Any rule that moves the origin cannot be linear.

\ex[Projection is linear]{Let $T:\RR^2\to\RR^2$ send each vector $v$ to its perpendicular projection $T(v)$ onto a fixed line through the origin. Geometrically this drops $v$ straight down onto the line. It is linear: projecting the sum $v+w$ gives the same point as projecting $v$ and $w$ separately and adding, and scaling $v$ scales its shadow. We have not chosen any basis, written any matrix, or computed anything --- yet we already know $T$ is linear and that $T(0)=0$.}

\ex[Rotation is linear]{Let $T:\RR^2\to\RR^2$ rotate every vector counterclockwise by $45^\circ$ about the origin. Rotating a parallelogram rotates its diagonal, so $T(v+w)=T(v)+T(w)$; rotating a stretched vector stretches the rotated vector, so $T(cv)=cT(v)$. Again the rule is described, and seen to be linear, with no coordinates at all. This is the picture Strang likes: ask what happens to a drawing of a house in the plane, and the answer (a tilted house) is obvious from the geometry, while it would be hard to guess from a matrix of numbers.}

\begin{center}
\begin{tikzpicture}[scale=0.95,>=Stealth]
  % house corners (original): base square (2,0)-(3.6,0)-(3.6,1.6)-(2,1.6),
  % roof apex (2.8,2.5), door (2.55,0)-(3.05,0.8)
  \pgfmathsetmacro{\axR}{(2.8-2.5)/sqrt(2)}
  \pgfmathsetmacro{\ayR}{(2.8+2.5)/sqrt(2)}

  % ---- light coordinate axes ----
  \draw[->,gray] (-3.2,0)--(4.4,0) node[right,font=\scriptsize]{$x$};
  \draw[->,gray] (0,-1.0)--(0,4.4) node[above,font=\scriptsize]{$y$};

  % ============ ORIGINAL HOUSE (gray) ============
  \begin{scope}[gray]
    \draw[very thick] (2,0)--(3.6,0)--(3.6,1.6)--(2,1.6)--cycle; % walls
    \draw[very thick] (2,1.6)--(2.8,2.5)--(3.6,1.6);            % roof
    \draw[thick] (2.55,0)--(2.55,0.8)--(3.05,0.8)--(3.05,0);    % door
  \end{scope}

  % ============ ROTATED HOUSE (blue) — image under T ============
  \begin{scope}[rotate around={45:(0,0)},blue!70!black]
    \draw[very thick] (2,0)--(3.6,0)--(3.6,1.6)--(2,1.6)--cycle;
    \draw[very thick] (2,1.6)--(2.8,2.5)--(3.6,1.6);
    \draw[thick] (2.55,0)--(2.55,0.8)--(3.05,0.8)--(3.05,0);
  \end{scope}

  % ---- trace the roof apex from original to rotated position ----
  \coordinate (Ap) at (2.8,2.5);
  \coordinate (ApR) at (\axR,\ayR);
  \fill[gray] (Ap) circle (1.6pt);
  \fill[blue!70!black] (ApR) circle (1.6pt);
  \pgfmathsetmacro{\rad}{sqrt(2.8*2.8+2.5*2.5)}
  \pgfmathsetmacro{\angA}{atan2(2.5,2.8)}
  \draw[->,thin,gray,dashed] (Ap) arc[start angle=\angA, end angle=\angA+45, radius=\rad];
  \node[font=\scriptsize,gray] at (2.05,3.55) {$45^\circ$};

  % labels
  \node[gray,font=\scriptsize] at (2.8,-0.55) {house};
  \node[blue!70!black,font=\scriptsize] at (-1.95,1.55) {$T(\text{house})$};
\end{tikzpicture}
\end{center}

It is just as important to see what is \emph{not} linear, because the failures are instructive.

\ex[Two non-examples]{The \emph{shift} $T(v)=v+v_0$ that slides the whole plane by a fixed nonzero vector $v_0$ is not linear: it moves the origin, $T(0)=v_0\neq0$. Concretely $T(2v)=2v+v_0$ while $2T(v)=2v+2v_0$, and these differ. (A map of this affine form is useful in graphics, but it is not a linear transformation.) The \emph{length} map $T(v)=\norm{v}$ also fails: it sends a vector to a single number, and $T(cv)=\abs{c}\norm{v}$, which is not $cT(v)$ when $c<0$. From here on $T$ always means a linear transformation.}

\ex[Multiplication by a matrix]{Every matrix $A\in\RR^{m\times n}$ \emph{defines} a linear transformation $T(v)=Av$, because matrix multiplication distributes over addition and pulls out scalars:
\[
   A(v+w)=Av+Aw, \qquad A(cv)=c\,Av .
\]
For instance $A=\left[\begin{smallmatrix}1&0\\0&-1\end{smallmatrix}\right]$ leaves the $x$-component of a vector alone and flips the sign of the $y$-component --- it reflects the plane across the $x$-axis. A $2\times3$ matrix gives a linear map $\RR^3\to\RR^2$, squashing space down a dimension. The plan for the rest of the section is the converse: \emph{every} linear transformation, once we pick a basis, is multiplication by some matrix.}

\subsection{A transformation is determined by its action on a basis}

How much do we need to know about $T$ to know it completely --- to predict $T(v)$ for \emph{every} input $v$? Far less than you might fear. Linearity is a powerful constraint.

Suppose we know $T(v_1)$ for one vector $v_1$. Then we know $T(cv_1)=c\,T(v_1)$ for every scalar multiple along that line. If we also know $T(v_2)$ for a second, independent vector, then for any combination we get
\[
   T(c_1 v_1 + c_2 v_2) = c_1 T(v_1) + c_2 T(v_2),
\]
so $T$ is determined on the entire plane they span. Push this to a full basis and the conclusion is complete.

\thm{A Transformation Is Determined by a Basis}{Let $v_1,\dots,v_n$ be a basis for the input space $\RR^n$, and let $T:\RR^n\to\RR^m$ be linear. Then $T$ is completely determined by the $n$ output vectors $T(v_1),\dots,T(v_n)$. Indeed, every input has a unique expansion $v=c_1 v_1+\cdots+c_n v_n$, and
\[
   T(v) = c_1\,T(v_1) + c_2\,T(v_2) + \cdots + c_n\,T(v_n).
\]}
\pf{Because $v_1,\dots,v_n$ is a basis, every $v\in\RR^n$ is a combination $v=c_1v_1+\cdots+c_nv_n$, and the coefficients $c_i$ are unique (Chapter~\ref{ch:4}: spanning gives existence, independence gives uniqueness). Applying linearity term by term,
\[
   T(v)=T\!\Big(\sum_i c_i v_i\Big)=\sum_i c_i\,T(v_i).
\]
The right-hand side uses only the known vectors $T(v_i)$ and the coordinates $c_i$ of $v$. So $T(v)$ is determined for every $v$.}

This theorem is the whole bridge from the coordinate-free world to matrices. The unique coefficients $c_1,\dots,c_n$ are the \emph{coordinates} of $v$ in the basis $v_1,\dots,v_n$. Coordinates come \emph{from} a basis; change the basis and the same vector gets new coordinates. The transformation itself does not change --- only its bookkeeping does.

\section{The Matrix of a Linear Transformation}

We now make the bridge explicit. To turn a linear transformation $T:\RR^n\to\RR^m$ into a matrix, we must choose \emph{two} bases: an input basis $v_1,\dots,v_n$ for $\RR^n$ to give coordinates to inputs, and an output basis $w_1,\dots,w_m$ for $\RR^m$ to give coordinates to outputs. The recipe for the matrix is forced on us by the theorem above.

\state{How to build the matrix of $T$}{Fix an input basis $v_1,\dots,v_n$ and an output basis $w_1,\dots,w_m$. To find column $j$ of the matrix $A$, apply $T$ to the $j$-th input basis vector and expand the result in the output basis:
\[
   T(v_j) = a_{1j}w_1 + a_{2j}w_2 + \cdots + a_{mj}w_m .
\]
The coefficients $a_{1j},\dots,a_{mj}$ are exactly column $j$ of $A$. With this $A$, if $v$ has coordinates $c=(c_1,\dots,c_n)$ in the input basis, then $T(v)$ has coordinates $Ac$ in the output basis.}

The rule looks circular at first, but it is the only thing it could be. Column $j$ records ``where the $j$-th input basis vector goes,'' written in the output basis. Linearity then propagates that information to every input, because $A c = \sum_j c_j(\text{column }j)$ reproduces $\sum_j c_j T(v_j) = T(v)$. The product $Ac$ being a combination of the columns --- the column picture from Chapter~\ref{ch:1} --- is exactly what makes this work.

\subsection{Projection: a clever basis versus the standard one}

The single best illustration is projection, because the right basis makes the matrix almost trivial and the wrong basis makes it messy --- for the very same transformation.

\ex[Projection onto a line, in two bases]{Let $T:\RR^2\to\RR^2$ project onto a line $L$ through the origin, and use the same basis for input and output ($n=m=2$).

\emph{Smart basis.} Choose $v_1=w_1$ to be a unit vector \emph{along} $L$ and $v_2=w_2$ a unit vector \emph{perpendicular} to $L$. Projection keeps the part along $L$ and kills the perpendicular part:
\[
   T(v_1)=v_1, \qquad T(v_2)=0 .
\]
Reading off coordinates in the basis $\{w_1,w_2\}$, the matrix is the diagonal
\[
   A=\begin{bmatrix} 1 & 0 \\ 0 & 0 \end{bmatrix},
   \qquad
   A\begin{bmatrix} c_1 \\ c_2\end{bmatrix}=\begin{bmatrix} c_1 \\ 0\end{bmatrix}.
\]
A perfect matrix: it simply zeroes the second coordinate. This is no accident --- $v_1,v_2$ are eigenvectors of the projection, with eigenvalues $1$ and $0$, and \textbf{in a basis of eigenvectors the matrix of any transformation is the diagonal matrix $\Lambda$ of eigenvalues} (Chapter~\ref{ch:7}).

\emph{Standard basis.} Now project onto the line at $45^\circ$ but insist on the standard basis $e_1=(1,0)$, $e_2=(0,1)$. With $a=(1,1)/\sqrt2$ a unit vector along the line, the projection matrix is $P=aa\T$ (Chapter~\ref{ch:5}),
\[
   P=\frac{a\,a\T}{a\T a}=\begin{bmatrix} \tfrac12 & \tfrac12 \\[2pt] \tfrac12 & \tfrac12\end{bmatrix}.
\]
Same transformation, different basis, and a fuller matrix. The lesson Strang draws is that the difficulty lives in the basis, not in the transformation. The diagonal $\left[\begin{smallmatrix}1&0\\0&0\end{smallmatrix}\right]$ and the dense $\left[\begin{smallmatrix}1/2&1/2\\1/2&1/2\end{smallmatrix}\right]$ describe one and the same projection.}

\subsection{The derivative is a matrix}

Linear transformations are not confined to arrows in $\RR^n$. The space of polynomials is a vector space (Chapter~\ref{ch:4}), and differentiation is linear on it --- a striking example because calculus turns into a matrix multiply.

\ex[The derivative as a matrix]{Let the input space be quadratic polynomials $c_1 + c_2 x + c_3 x^2$, with basis $v_1=1$, $v_2=x$, $v_3=x^2$. The derivative
\[
   T(c_1 + c_2 x + c_3 x^2) = c_2 + 2c_3 x
\]
lands in the space of linear polynomials, with output basis $w_1=1$, $w_2=x$. To build the matrix, differentiate each input basis vector and expand:
\[
   T(1)=0,\qquad T(x)=1=w_1,\qquad T(x^2)=2x=2w_2 .
\]
The three results, in output coordinates, are the columns of
\[
   A=\begin{bmatrix} 0 & 1 & 0 \\ 0 & 0 & 2 \end{bmatrix},
   \qquad
   A\begin{bmatrix} c_1 \\ c_2 \\ c_3\end{bmatrix}=\begin{bmatrix} c_2 \\ 2c_3\end{bmatrix}.
\]
Differentiating $c_1+c_2x+c_3x^2$ has become multiplying its coordinate vector by $A$, and the output coordinates $(c_2,2c_3)$ are exactly the coefficients of $c_2+2c_3x$. The matrix \emph{is} the derivative, once a basis is fixed.}

\rmkb[Where matrix multiplication comes from]{Two facts close the loop and explain operations we have used all along. If $T$ is invertible, then the inverse transformation $T^{-1}$ has matrix $A^{-1}$ --- undoing the map undoes the matrix. And if we do one transformation after another, $T_2\circ T_1$ (first $T_1:v\mapsto A_1v$, then $T_2:w\mapsto A_2w$), the combined transformation has matrix $A_2A_1$. \textbf{Matrix multiplication is defined the way it is precisely so that it matches composition of transformations.} The row-times-column rule is not arbitrary; it is forced by ``do $A_1$, then $A_2$.''}

\section{Change of Basis and Similar Matrices}

We have seen one transformation wear two different matrices, depending on the basis. The relationship between those matrices is not random --- it is a precise algebraic tie called \emph{similarity}, and it is the engine behind eigenvalues, diagonalization, and the compression scheme in the next section. First we settle how the \emph{coordinates of a vector} change, then how the \emph{matrix of a transformation} changes.

\subsection{Coordinates of a vector}

Suppose we install a new basis whose vectors are the columns of a matrix $W=\br{w_1\ w_2\ \cdots\ w_n}$. A vector $x$, written in the \emph{standard} basis, has new coordinates $c=(c_1,\dots,c_n)$ in the $w$-basis meaning $x=c_1w_1+\cdots+c_nw_n$. That sum is exactly $Wc$:

\state{Change of basis for vectors}{If the columns of $W$ are the new basis vectors, then standard coordinates $x$ and new coordinates $c$ are related by
\[
   x = Wc, \qquad\text{equivalently}\qquad c = W^{-1}x .
\]
To \emph{express} $x$ in the new basis we solve for its coordinates $c=W^{-1}x$; to \emph{reconstruct} $x$ from coordinates we multiply $x=Wc$.}

The matrix $W$ is invertible precisely because its columns are a basis (independent and spanning), so $W^{-1}$ exists and the dictionary between $x$ and $c$ goes both ways.

\subsection{The matrix of a transformation under change of basis}

Now the main event. One transformation $T$, two bases, two matrices --- how are they related? Suppose $T$ has matrix $A$ in the old basis and matrix $B$ in the new basis, with $M$ the change-of-basis matrix linking the two coordinate systems. Tracking a vector through both descriptions forces the relationship.

\thm{Change of Basis Produces Similar Matrices}{Let $T$ be a linear transformation with matrix $A$ in one basis and matrix $B$ in another. If $M$ is the change-of-basis matrix from the new coordinates to the old (so an input written $u$ in the new basis is $Mu$ in the old), then
\[
   B = M^{-1} A\, M .
\]
Two matrices related this way are called \textbf{similar}.}
\pf{Take an input with new coordinates $u$. Its old coordinates are $Mu$. Apply $T$ in old coordinates: the output, in old coordinates, is $A(Mu)$. Convert that output back to new coordinates by multiplying by $M^{-1}$: the output in new coordinates is $M^{-1}AMu$. But by definition $B$ is the matrix that sends new input coordinates straight to new output coordinates, so $Bu = M^{-1}AMu$ for every $u$. Hence $B=M^{-1}AM$.}

Similarity is the algebraic fingerprint of ``same transformation, different basis.'' Similar matrices share everything that belongs to the transformation rather than to the basis: the same eigenvalues, the same determinant, the same rank, the same trace, the same characteristic polynomial. They differ only in the coordinates used to write them down.

\ex[Diagonalization is a change of basis]{The eigenvalue factorization $A=S\Lambda S^{-1}$ from Chapter~\ref{ch:7} is exactly this theorem read backwards. Rearranged, $\Lambda = S^{-1}A S$, which is $B=M^{-1}AM$ with $M=S$ the eigenvector matrix and $B=\Lambda$. The message: \emph{in the basis of eigenvectors, the transformation's matrix is the diagonal $\Lambda$.} Diagonalizing a matrix means choosing the basis in which the transformation is simplest --- pure stretching along the eigenvector axes, with no mixing. When the basis is orthonormal eigenvectors (the symmetric case, Chapter~\ref{ch:8}), $M=Q$ is orthogonal and $M^{-1}=Q\T$, so the change of basis costs only a transpose.}

\section{Application: Image Compression}

Here is why a good basis is worth real money. The whole content of compression is: \emph{store the same information in a basis where most coordinates are nearly zero, then throw the small ones away.}

Picture one black-and-white video frame at $512\times 512$ pixels. The camera records a brightness level for each of the $512^2$ pixels, so a single frame is a vector in a $512^2$-dimensional space. The \emph{standard basis} for that space has one vector per pixel, and transmitting all $512^2$ coordinates frame after frame would swamp any channel. The trouble is that the standard basis is blind to structure: there is no efficient way for it to say ``the whole frame is black,'' because that statement still costs one number per pixel.

Change the basis and the costs change. Suppose one basis vector is the all-ones vector $(1,1,\dots,1)$ --- every pixel at the same brightness. Then a blank gray board is a \emph{single} coordinate, not $512^2$ of them. Add a vector that is half $+1$ and half $-1$ (bright on the left, dark on the right), a vector that alternates $+1,-1,+1,-1$ (fine stripes), and so on for the various patterns an image is built from. Smooth, large-scale features now live in just a few coordinates.

\ex[Two good bases: Fourier and Haar]{The best-known choice is the \textbf{Fourier basis}, whose vectors are sampled cosines and sines (the real parts of the powers $\omega^{jk}$ from the Fourier matrix). The JPEG standard cuts each frame into $8\times 8$ blocks of $64$ pixels and rewrites every block in this cosine basis. A competing choice is the \textbf{Haar wavelet basis} for $\RR^8$,
\[
   \begin{bmatrix} 1\\1\\1\\1\\1\\1\\1\\1\end{bmatrix},\
   \begin{bmatrix} 1\\1\\1\\1\\-1\\-1\\-1\\-1\end{bmatrix},\
   \begin{bmatrix} 1\\1\\-1\\-1\\0\\0\\0\\0\end{bmatrix},\
   \begin{bmatrix} 0\\0\\0\\0\\1\\1\\-1\\-1\end{bmatrix},\ \dots
\]
The first vector measures overall brightness; later vectors measure finer and finer differences, many of their entries being $0$. These vectors are orthogonal and can be scaled to orthonormal, which is what makes the next step cheap. JPEG2000 uses an improved wavelet basis of this kind.}

The mechanics are nothing but change of basis. To rewrite a signal $x$ (standard, pixel coordinates) in a new basis with vectors as the columns of $W$, we want coordinates $c$ with
\[
   x = c_1 w_1 + \cdots + c_n w_n = Wc,
   \qquad\text{so}\qquad
   c = W^{-1}x .
\]
This is the change-of-basis formula again. Two properties make a basis good for compression:
\begin{itemize}
   \item \textbf{Fast transforms.} Multiplying by $W$ and by $W^{-1}$ must be cheap. The Fourier basis is ideal because the FFT computes $Wc$ and $W^{-1}x$ in $O(n\log n)$ time, and a wavelet transform is even faster. An \emph{orthonormal} basis is best of all: then $W^{-1}=W\T$, so the inverse transform is free.
   \item \textbf{Sparse coordinates.} Most of the $c_i$ should come out small, safely set to zero without a visible change. Storing only the few large coefficients is the compression.
\end{itemize}
The pipeline is: take $x$, compute its $64$ coefficients $c$ losslessly in the new basis, then \emph{lossily} discard the coefficients below a threshold, leaving $\hat c$ with many zeros, and reconstruct $\hat x=\sum \hat c_i w_i\approx x$.

\rmkb[Why not the eigenvector basis?]{In principle the perfect basis is the eigenvectors of the relevant transformation --- it diagonalizes everything. In practice computing that basis is far more expensive than a JPEG could ever afford. Fourier and wavelet bases are the engineering compromise: not optimal, but \emph{fixed} (the same for every image) and equipped with a fast transform. The best basis for a slide of algebra notes differs from the best basis for an action movie, but both pay off for the same reason --- they concentrate the picture into a handful of coordinates.}

\section{Left and Right Inverses}

We return to inverses, but now for matrices that are not square --- or that are square but singular. A genuine, two-sided inverse $A^{-1}$ satisfies $AA^{-1}=I=A^{-1}A$ and exists only when $A$ is square with full rank, $r=m=n$. Then $N(A)=\{0\}$ and $N(A\T)=\{0\}$: nothing is collapsed in either direction. When $A$ is rectangular, we cannot have a two-sided inverse, because either $A$ or $A\T$ has a nontrivial null space. But we can sometimes invert from \emph{one} side, and which side depends on which kind of full rank $A$ has.

\subsection{Full column rank: a left inverse}

Suppose $A$ is $m\times n$ with \emph{full column rank} $r=n$ (tall, independent columns). Then $N(A)=\{0\}$: the equation $Ax=b$ has \emph{at most} one solution --- exactly one when $b\in C(A)$, none otherwise. From Chapter~\ref{ch:5} we know the crucial fact that makes least squares work: when the columns of $A$ are independent, the $n\times n$ matrix $A\T A$ is invertible.

\state{The left inverse}{If $A$ has full column rank $r=n$, then $A\T A$ is invertible and
\[
   A^{-1}_{\text{left}} = (A\T A)^{-1}A\T
\]
is a \textbf{left inverse}: it satisfies $A^{-1}_{\text{left}}\,A = (A\T A)^{-1}A\T A = I_n$. Multiplying $A$ on the \emph{left} by this matrix recovers the identity.}

This is precisely the matrix that solved least squares: $\hat x=(A\T A)^{-1}A\T b$ is the best solution of an unsolvable $Ax=b$. The left inverse undoes $A$ from the left, but the \emph{other} product,
\[
   A\,A^{-1}_{\text{left}} = A(A\T A)^{-1}A\T = P,
\]
is the projection matrix onto the column space $C(A)$ (Chapter~\ref{ch:5}), not the identity --- it equals $I_m$ only when $m=n$. A tall matrix can be undone after it acts, but it cannot be undone before: there is no way to recover the directions of $\RR^m$ outside the column space, because $A$ never reached them.

\subsection{Full row rank: a right inverse}

The mirror image. Suppose $A$ is $m\times n$ with \emph{full row rank} $r=m$ (wide, independent rows). Then $N(A\T)=\{0\}$ and $C(A)=\RR^m$, so $Ax=b$ is solvable for \emph{every} $b$ --- in fact, with $n-m$ free variables, there are infinitely many solutions when $n>m$. Now $AA\T$ is the invertible one (it is $m\times m$ and full rank).

\state{The right inverse}{If $A$ has full row rank $r=m$, then $AA\T$ is invertible and
\[
   A^{-1}_{\text{right}} = A\T(AA\T)^{-1}
\]
is a \textbf{right inverse}: it satisfies $A\,A^{-1}_{\text{right}} = AA\T(AA\T)^{-1} = I_m$. Multiplying $A$ on the \emph{right} by this matrix recovers the identity.}

Symmetrically, the opposite product
\[
   A^{-1}_{\text{right}}\,A = A\T(AA\T)^{-1}A
\]
is the projection onto the row space $C(A\T)$, not the identity. The right inverse picks out, among the infinitely many solutions of $Ax=b$, the one lying in the row space.

\rmkb[One table to remember]{The two one-sided inverses are transposes of each other's construction, and each lives where its $A\T A$ or $AA\T$ is invertible.
\begin{center}
\begin{tabular}{lll}
\toprule
case & condition & inverse \\
\midrule
two-sided & $r=m=n$ & $A^{-1}$ \\
left & $r=n<m$ (full column rank) & $(A\T A)^{-1}A\T$ \\
right & $r=m<n$ (full row rank) & $A\T(AA\T)^{-1}$ \\
\bottomrule
\end{tabular}
\end{center}
The remaining and hardest case --- $r<m$ \emph{and} $r<n$, short rank on both sides --- has neither a left nor a right inverse. That is the case the pseudoinverse is built for.}

\section{The Pseudoinverse \texorpdfstring{$A^{+}$}{A+}}

When $A$ has $r<m$ and $r<n$, both $N(A)$ and $N(A\T)$ are nontrivial. A nonzero $x$ with $Ax=0$ kills information that no inverse can resurrect: there is no matrix sending $0$ back to $x$. So we cannot ask for a true inverse. But there is a clean part of the action that \emph{is} invertible, and the pseudoinverse inverts exactly that.

The clean part is the heart of the four-subspaces picture. Recall (Chapter~\ref{ch:4}) that $A$ crushes the null space $N(A)$ to zero, and on the row space $C(A\T)$ --- the orthogonal complement of $N(A)$ inside $\RR^n$ --- it acts one-to-one onto the column space $C(A)$. Both of those spaces have dimension $r$. So restricted to the row space, $A$ is a genuine bijection onto the column space, and \emph{that} restriction can be inverted.

\thm{Row Space and Column Space Are in One-to-One Correspondence}{Let $A$ be $m\times n$ of rank $r$. The map $x\mapsto Ax$ sends the row space $C(A\T)$ one-to-one onto the column space $C(A)$. That is, if $x\neq y$ are both in the row space, then $Ax\neq Ay$.}
\pf{Suppose instead that $x\neq y$ lie in the row space but $Ax=Ay$. Then $A(x-y)=0$, so $x-y\in N(A)$. But the row space is a subspace, so the difference $x-y$ of two row-space vectors is again in the row space. Thus $x-y$ lies in both $N(A)$ and $C(A\T)$. By the Fundamental Theorem these are orthogonal complements (Chapter~\ref{ch:4}), meeting only at the origin, so $x-y=0$, i.e.\ $x=y$ --- contradicting $x\neq y$. Hence distinct row-space vectors have distinct images: the map is one-to-one, and since its image is all of $C(A)$ (every $Ax$ lies in $C(A)$, and the row-space part of any $x$ already produces it), it is onto.}

\begin{center}
\begin{tikzpicture}[scale=1.0,>=Stealth,font=\footnotesize]
  % ---------- geometry constants ----------
  \def\bw{4.6}   % box width
  \def\bh{6.2}   % box height
  \def\gap{3.4}  % horizontal gap between boxes
  % wedge geometry (symmetric about horizontal axis through the origin)
  \def\L{2.55}
  \def\cx{2.311}  % \L*cos(25)
  \def\cy{1.078}  % \L*sin(25)

  \coordinate (Lo) at (0,0);
  \coordinate (Ro) at ({\bw+\gap},0);

  \draw[rounded corners=10pt,thick,gray!55] (Lo) rectangle ++(\bw,\bh);
  \draw[rounded corners=10pt,thick,gray!55] (Ro) rectangle ++(\bw,\bh);

  % ===================== LEFT: input space R^n =====================
  \coordinate (OL) at ($(Lo)+(1.05,\bh*0.5)$);
  % row space wedge (blue, up-right), null space wedge (red, down-right): perpendicular
  \fill[blue!10]  (OL) -- ($(OL)+(1.95,1.30)$) -- ($(OL)+(2.55,0.62)$) -- cycle;
  \fill[red!10]   (OL) -- ($(OL)+(2.55,-0.62)$) -- ($(OL)+(1.95,-1.30)$) -- cycle;
  \draw[very thick,blue!70!black] (OL) -- ($(OL)+(\cx,\cy)$);
  \draw[very thick,red!70!black]  (OL) -- ($(OL)+(\cx,-\cy)$);
  % right-angle marker at OL
  \draw[gray] ($(OL)+(0.544,0.254)$) -- ($(OL)+(0.811,0)$) -- ($(OL)+(0.544,-0.254)$);

  % labels left
  \node[blue!70!black,align=center] at ($(OL)+(2.35,1.95)$) {row space\\ $\Col{A\T}$\\ $\dim = r$};
  \node[red!70!black,align=center]  at ($(OL)+(2.35,-1.95)$) {null space\\ $\Nul{A}$\\ $\dim = n-r$};
  \node[font=\scriptsize,gray!55!black] at ($(Lo)+(\bw*0.5,\bh+0.30)$) {\itshape input space $\RR^n$};

  % representative vector in the row space
  \coordinate (XR) at ($(OL)+(1.50,0.70)$);
  \fill[blue!70!black] (XR) circle (1.7pt);
  \node[blue!70!black,below=1pt,font=\scriptsize] at (XR) {$x$};

  \fill (OL) circle (1.5pt);
  \node[left,font=\scriptsize] at (OL) {$0$};

  % ===================== RIGHT: output space R^m =====================
  \coordinate (OR) at ($(Ro)+(1.05,\bh*0.5)$);
  \fill[blue!10]  (OR) -- ($(OR)+(1.95,1.30)$) -- ($(OR)+(2.55,0.62)$) -- cycle;
  \fill[red!10]   (OR) -- ($(OR)+(2.55,-0.62)$) -- ($(OR)+(1.95,-1.30)$) -- cycle;
  \draw[very thick,blue!70!black] (OR) -- ($(OR)+(\cx,\cy)$);
  \draw[very thick,red!70!black]  (OR) -- ($(OR)+(\cx,-\cy)$);
  \draw[gray] ($(OR)+(0.544,0.254)$) -- ($(OR)+(0.811,0)$) -- ($(OR)+(0.544,-0.254)$);

  \node[blue!70!black,align=center] at ($(OR)+(2.35,1.95)$) {column space\\ $\Col{A}$\\ $\dim = r$};
  \node[red!70!black,align=center]  at ($(OR)+(2.35,-1.95)$) {left null space\\ $\Nul{A\T}$\\ $\dim = m-r$};
  \node[font=\scriptsize,gray!55!black] at ($(Ro)+(\bw*0.5,\bh+0.30)$) {\itshape output space $\RR^m$};

  % image b = A x  in the column space
  \coordinate (B) at ($(OR)+(1.50,0.70)$);
  \fill[blue!70!black] (B) circle (1.7pt);
  \node[blue!70!black,below right,font=\scriptsize] at (B) {$Ax$};

  \fill (OR) circle (1.5pt);
  \node[right,font=\scriptsize] at (OR) {$0$};

  % ===================== the bijection: two opposing arrows =====================
  % forward A : row space -> column space (one-to-one onto), arcs over the top
  \draw[->,thick,blue!70!black] (XR) to[bend left=22] (B);
  \node[blue!70!black,font=\scriptsize,fill=white,inner sep=1.5pt]
       at ($(XR)!0.5!(B)+(0,1.30)$) {$A$ \ (one-to-one onto)};

  % reverse A^+ : column space -> row space, arcs under the forward arrow
  \draw[->,thick,blue!70!black,dashed] (B) to[bend left=22] (XR);
  \node[blue!70!black,font=\scriptsize,fill=white,inner sep=1.5pt]
       at ($(XR)!0.5!(B)+(0,0.30)$) {$A^{+}$ \ (inverts it)};

  % ===================== the collapses to 0 =====================
  % null space N(A) -> 0 in the output space (A crushes it)
  \coordinate (XN) at ($(OL)+(1.50,-0.70)$);
  \fill[red!70!black] (XN) circle (1.7pt);
  \draw[->,thick,red!70!black,dashed] (XN) to[bend left=12] (OR);
  \node[red!70!black,font=\scriptsize,align=center,fill=white,inner sep=1.5pt]
       at ($(XN)+(1.55,0.46)$) {$A$ crushes\\ $\Nul{A}\to 0$};

  % left null space N(A^T) -> 0 in the input space (A^+ sends it to 0)
  \coordinate (YN) at ($(OR)+(1.50,-0.70)$);
  \fill[red!70!black] (YN) circle (1.7pt);
  \draw[->,thick,red!70!black,dashed] (YN) to[bend left=12] (OL);
  \node[red!70!black,font=\scriptsize,align=center,fill=white,inner sep=1.5pt]
       at ($(YN)+(-1.55,-0.46)$) {$A^{+}$ sends\\ $\Nul{A\T}\to 0$};
\end{tikzpicture}
\end{center}

This bijection is what the pseudoinverse inverts. It returns each $Ax$ in the column space to the unique $x$ in the row space that produced it, and it sends everything in the left null space $N(A\T)$ to zero (those directions never came from any input we can recover).

\defn{Pseudoinverse}{The \textbf{pseudoinverse} $A^{+}$ (also called the Moore--Penrose inverse) of an $m\times n$ matrix $A$ is the $n\times m$ matrix that inverts the row-space-to-column-space bijection: $A^{+}(Ax)=x$ for every $x$ in the row space $C(A\T)$, and $A^{+}y=0$ for every $y$ in the left null space $N(A\T)$. In symbols, $N(A^{+})=N(A\T)$.}

\subsection{Building \texorpdfstring{$A^{+}$}{A+} from the SVD}

The definition is geometric; the construction is the singular value decomposition (Chapter~\ref{ch:9}). The SVD is the perfect tool because it presents $A$ already split along the four subspaces.

Write $A=U\Sigma V\T$, where $U$ ($m\times m$) and $V$ ($n\times n$) are orthogonal and $\Sigma$ ($m\times n$) holds the singular values $\sigma_1\ge\cdots\ge\sigma_r>0$ on its first $r$ diagonal entries, zeros elsewhere. The first $r$ columns of $V$ are an orthonormal basis for the row space, the first $r$ columns of $U$ an orthonormal basis for the column space, and $A$ stretches the $k$-th row-space axis by $\sigma_k$ onto the $k$-th column-space axis.

Because $U$ and $V$ are orthogonal, $U^{-1}=U\T$ and $V^{-1}=V\T$, so inverting $A$ reduces to inverting $\Sigma$. And $\Sigma$ is ``diagonal,'' so we invert it the only sensible way: reciprocate the nonzero singular values and leave the zeros alone.

\state{The pseudoinverse via SVD}{If $A=U\Sigma V\T$, then
\[
   A^{+} = V\,\Sigma^{+}\,U\T,
\]
where $\Sigma^{+}$ is the $n\times m$ matrix with $1/\sigma_1,\dots,1/\sigma_r$ on its first $r$ diagonal entries and zeros elsewhere. The product $\Sigma\Sigma^{+}$ is the $m\times m$ diagonal matrix with $r$ ones (then zeros); $\Sigma^{+}\Sigma$ is the $n\times n$ version. When $A$ is square and invertible ($r=m=n$), every $\sigma_i\neq0$, $\Sigma^{+}=\Sigma^{-1}$, and $A^{+}=A^{-1}$ collapses to the ordinary inverse.}

The reason we reciprocate $\sigma_k$ but leave the zeros as zeros is the geometry: on the row space, $A$ multiplies the $k$-th axis by $\sigma_k$, so undoing it divides by $\sigma_k$; in the null-space directions, $A$ multiplied by $0$, and there is nothing to undo, so $A^{+}$ multiplies by $0$ there too. That is the only choice consistent with $N(A^{+})=N(A\T)$.

\ex[A rank-one pseudoinverse]{Take the rank-one $A=\left[\begin{smallmatrix}1&2\\2&4\end{smallmatrix}\right]$, both columns multiples of $(1,2)$, so $r=1$, $m=n=2$. A two-sided inverse is impossible (the matrix is singular), yet $A^{+}$ exists. Writing $A=\sigma_1 u_1 v_1\T$ with unit vectors $u_1=v_1=(1,2)/\sqrt5$ and $\sigma_1=5$ (since $\norm{(1,2)}^2=5$ and $A=(1,2)(1,2)\T$), the pseudoinverse reciprocates the one nonzero singular value:
\[
   A^{+}=\tfrac{1}{\sigma_1}\,v_1 u_1\T
        =\tfrac{1}{5}\cdot\frac{1}{5}\begin{bmatrix}1\\2\end{bmatrix}\begin{bmatrix}1&2\end{bmatrix}
        =\frac{1}{25}\begin{bmatrix}1&2\\2&4\end{bmatrix}.
\]
Check the defining property on the row space: the row space is the line through $v_1$, and $A^{+}A$ projects onto it. Indeed $A^{+}A=\tfrac1{25}\left[\begin{smallmatrix}1&2\\2&4\end{smallmatrix}\right]\left[\begin{smallmatrix}1&2\\2&4\end{smallmatrix}\right]=\tfrac1{25}\left[\begin{smallmatrix}5&10\\10&20\end{smallmatrix}\right]=\tfrac15\left[\begin{smallmatrix}1&2\\2&4\end{smallmatrix}\right]$, the projection onto the line through $(1,2)$, exactly as the theory predicts.}

\subsection{The pseudoinverse as the shortest least-squares solution}

The pseudoinverse is not a curiosity. It is the single formula that solves $Ax=b$ ``as well as possible'' in \emph{every} case at once --- the unifying answer that ordinary, left, and right inverses each handled only in their special situation.

\state{What $A^{+}b$ computes}{For any $A$ and any $b$, the vector
\[
   \hat x = A^{+}b
\]
is the \textbf{minimum-norm least-squares solution} of $Ax=b$:
\begin{itemize}
   \item among all $x$ that minimize $\norm{Ax-b}$ (the least-squares condition), $\hat x$ is the one of smallest length $\norm{x}$;
   \item it always lies in the row space $C(A\T)$, the part of $\RR^n$ with no wasted null-space component.
\end{itemize}}

Each earlier inverse is a special case of this one. When $A$ is square and invertible, $A^{+}=A^{-1}$ and $\hat x$ is the exact solution. When $A$ has full column rank, $A^{+}=(A\T A)^{-1}A\T=A^{-1}_{\text{left}}$, and $\hat x$ is the least-squares solution of Chapter~\ref{ch:5}. When $A$ has full row rank, $A^{+}=A\T(AA\T)^{-1}=A^{-1}_{\text{right}}$, and $\hat x$ is the shortest of the infinitely many exact solutions. In the hardest case --- short rank on both sides, no left or right inverse --- $A^{+}b$ still returns the shortest best-fit. This is why statisticians reach for $A^{+}$ in regression: they cannot always promise their design matrix has independent columns, but $A^{+}$ returns a sensible answer regardless.

\rmkb[Everything ties together here]{Look at what the pseudoinverse gathers up. The \emph{four fundamental subspaces} (Chapter~\ref{ch:4}) provide the stage: $A^{+}$ inverts the row-space-to-column-space bijection and zeroes the left null space. The \emph{SVD} (Chapter~\ref{ch:9}) provides the construction: orthogonal $U,V$ and a reciprocated $\Sigma$. \emph{Projection and least squares} (Chapter~\ref{ch:5}) provide the meaning: $AA^{+}$ projects onto $C(A)$, $A^{+}A$ projects onto $C(A\T)$, and $A^{+}b$ is the shortest best-fit solution. The pseudoinverse is the place where the four subspaces, the SVD, orthogonal matrices, and least squares are revealed to be one circle of ideas --- the right note on which to close the book.}

\section{What to Carry Forward}

A matrix was never the primary object; it is the coordinate shadow of a \emph{linear transformation}, a coordinate-free rule satisfying $T(v+w)=Tv+Tw$ and $T(cv)=cTv$. Choosing an input basis and an output basis turns $T$ into a matrix $A$ whose $j$-th column is $T(v_j)$ written in the output basis. Choose the basis well --- eigenvectors, Fourier vectors, wavelets --- and the matrix becomes simple (diagonal, or fast to apply); that single fact powers diagonalization and image compression alike.

Because the matrix depends on the basis, one transformation has many matrices, all linked by $B=M^{-1}AM$: they are \emph{similar}, sharing eigenvalues, determinant, rank, and trace while differing only in coordinates. And when a matrix has no honest inverse, the rank decides the substitute: full column rank gives the left inverse $(A\T A)^{-1}A\T$, full row rank gives the right inverse $A\T(AA\T)^{-1}$, and the general case gives the pseudoinverse $A^{+}=V\Sigma^{+}U\T$, which inverts $A$ on the row space, kills the left null space, and returns the shortest least-squares solution $\hat x=A^{+}b$ for any $A$ and any $b$. Four subspaces, one number $r$, and the SVD --- the threads of the whole book pulled into a single knot.
